\subsubsection{Sample Preparation -- Cell lysis and Protein digestion} %\label{sec:methods-ont-prep}

Cell pellets of~\syn~and~\syna~strains were mixed with 100 $\mu$L of SDS-based lysis buffer (4\% SDS, 100 mM Tris-HCl, pH 8.0) containing final concentration of 10 mM Tris (2-carboxyethyl) phosphine (TCEP) and 40 mM chloroacetamide (CAA), vortexed for 5 min, and then boiled at 95\textdegree C for 10 min.
The lysate samples were then processed for proteomics analysis following the E3technology procedure reported recently \citep{martin_development_2024}.
In brief, the lysates were mixed with 4$\times$ volume of 80\% acetonitrile (ACN) and then transferred to E3filters (CDS Analytical, Oxford, PA) followed by centrifugation at 3,000 rpm for 2 min.
The filters were washed with 500 $\mu$L of 80\% ACN for three times, and then transferred to clean collection tubes.
For protein digestion, around 1.0 $\mu$g of trypsin and 200 $\mu$L of digestion buffer (50 mM triethylammonium bicarbonate, TEAB) were added to each sample followed by incubation at 37\textdegree C for 16-18 hours with gentle shaking (300 rpm/min).
After digestion, the filters were first spun at 3,000 rpm for 2 min to collect the flow through, and then eluted sequentially with 200 $\mu$L of 0.2\% formic acid (FA) in water and 200 $\mu$L of 0.2\% FA in 50\% ACN.
The pooled elution was dried in SpeedVac, and stored under -80\textdegree C until further analysis.

\subsubsection{LC-MS/MS analysis and protein identification}
The LC-MS/MS analysis was performed using an Ultimate 3000 RSLCnano system in conjunction with an Orbitrap Eclipse mass spectrometer and FAIMS Pro Interface (Thermo Scientific), similar to the procedures described recently \citep{sun_-cell_2026}.
In brief, the peptides resuspended in LC buffer A (0.1\% FA in water) were first loaded onto a trap column (PepMap100 C18,300 $\mu$m $\times$ 2 mm, 5 $\mu$m particle, Thermo Scientific), and then separated on an analytical column (PepMap100 C18, 50 cm $\times$ 75 $\mu$m i.d., 3 $\mu$m; Thermo Scientific) at a flow of 250 nL/min.
A linear LC gradient was applied from 1\% to 25\% mobile phase B (0.1\% formic acid in acetonitrile) over 125 min, followed by an increase to 32\% mobile phase B over 10 min.
The column was washed with 80\% mobile phase B for 5 min, followed by equilibration with mobile phase A for 15 min.
For MS analysis, the Detector type was Orbitrap with resolution of 120,000; Precursor MS range (m/z) was 380–980; AGC target was Standard; Maximum injection time mode was Auto.
For data-independent acquisition (DIA) MS/MS analysis, the Isolation mode was Quadruple; DIA Window type was Auto and Isolation Window (m/z) was 8 with an overlap of 1; activation type was HCD with fixed collision energy mode (30\%); the Detector Type was Orbitrap with a resolution of 30,000; Normalized AGC target (\%) was 800, and the Maximum injection time mode was Auto; the Loop Control was 2 sec.
For FAIMS compensation voltages (CV) setting, a 3-CV combination ($\sim$40, $\sim$55, and $\sim$75) was applied.

For proteome identification, the MS raw data were processed using Spectronaut software (version 19), a library-free DIA analysis workflow with directDIA+ and the in-house curated protein databases (Syn1.0, 828 sequences; Syn3A, 458 sequences).
Detailed parameters for Pulsar and library generation include: Trypsin/P as specific enzyme with minimum peptide length of 7 amino acids and maximally 2 missed cleavages; toggle N-terminal M turned on; Oxidation on M and Acetyl at protein N-terminus as variable modifications; Carbamidomethyl on C as fixed modification; FDRs at PSM, peptide and protein level all set to 0.01; Quantity MS level set to MS2, and cross-run normalization turned on.
The MS raw data associated with this study has been deposited onto the MassIVE Repository with the accession MSV000099558 (password: 4zzHbHHqqW7Gl2F5).

\subsubsection{Relative Protein Quantification}

Protein abundances were quantified from the Spectronaut iBAQ (intensity-based absolute quantification) values, which serve as a within-sample proxy for molar protein amount.
For each of the three replicates per organism, iBAQ was converted to a relative abundance in parts per million (iPM), $\text{iPM}_i = 10^{6}\,\text{iBAQ}_i / \sum_j \text{iBAQ}_j$, and the per-protein mean iPM across replicates was used for all downstream analyses, with the across-replicate coefficient of variation retained as a quality metric.
For~\syn, 721 of the 828 annotated protein-coding genes (87.1\%) carried a quantified iPM.
The per-gene relative abundances (iPM) are available in \sdomics.

\subsubsection{Absolute Protein Quantification}

%Currently in \refsec{methods-ptn-quant}. \cmt{Do we want to move this here?}
%\subsubsection{Proteome Quantification for \jsyn}
\label{sec:methods-ptn-quant}

The~\syn~mass spectrometry data provides relative abundances for 721 of the 828 protein-coding genes, which can be converted to absolute abundances by scaling the relative amount with the total number of proteins in the cell.
The total number of proteins has not been measured directly within \syn, but can be estimated via an approach similar to that used in \cite{breuer_essential_2019}.
First the dry mass of the cell and protein mass fraction are needed.
The protein mass fraction is assumed to be 58.2\%, which was measured in the parent organism, \Mmy~\citep{razin_chemical_1963}.
The dry mass of the cell, 12.8 fg, was also estimated identically to the dry mass of \jsyna~(see \cite{breuer_essential_2019}, \textit{Appendix 1}).
Accordingly, the average \jsyn~cell has approximately 127,000 proteins; the cell volume is calculated as the sphere of 220 nm radius as 0.044 $\mu m^3$ \cite{moger-reischer_evolution_2023}.
Thus, a protein density of $2.87\times 10^{6}$ $proteins/\mu m^{3}$ is estimated, which is in rough agreement with $3.5-4.4 \times 10^{6}$ $proteins/\mu m^{3}$ measured in \eco~\citep{milo_what_2013}.
Finally, absolute protein abundances were calculated using the following formula:
\begin{equation}
\label{eqn:abs-formula}
\text{$Abs~Abundance_{i} = \frac{dry~mass * mass~fraction}{avg~molecular~weight} * Rel~Abundance_{i}$}
\end{equation}
The resulting per-gene absolute copy numbers are available in \sdomics.

\subsubsection{Localization of Proteome}

Subcellular localization was predicted from sequence for every~\syn~protein, since membrane-associated proteins are systematically under-recovered by the trypsin-digestion workflow and are therefore handled separately in the transcriptome--proteome correlation.
Transmembrane regions (TMRs) were predicted with \texttt{DeepTMHMM}~\cite{hallgren_deeptmhmm_2022} and signal peptides with \texttt{SignalP} 6.0~\cite{teufel_signalp_2022}.
Each protein was assigned to one of four classes by priority: a SignalP secretory or lipoprotein signal peptide marked it as extracellular or lipoprotein, respectively; otherwise one or more DeepTMHMM TMRs marked it as membrane, and proteins with no TMR were assigned cytoplasmic.
Of the 721 detected~\syn~proteins, 516 were cytoplasmic, 126 membrane, 68 lipoprotein, and 11 extracellular.
Median copy number per cell fell steeply with membrane association, from 47 for cytoplasmic proteins to 21 (lipoprotein), 10 (membrane), and 3 (extracellular), consistent with the under-digestion and under-recovery of membrane-embedded proteins rather than a true abundance difference.
The per-protein localization assignments are available in \sdomics.

